\section{Discussion and conclusion}

The most useful interpretation of the Genesis IV result is not that a software organism became broadly more intelligent twice. It is that an improvement mechanism became measurably more selective about when expensive search was worth continuing, and that this selectivity itself survived a second recursive transition.

The path to that result was not monotone. A bounded parameter family first produced G1. Attempts to recurse within that family saturated. Expanding to executable policy code then produced an attractive but brittle shortcut that failed every fresh V18 holdout comparison by tying the parent. Only after the development distribution was diversified did a G2 policy emerge that generalized prospectively. That history suggests a practical lesson shared with recent work on harness regularization: recursive systems need explicit mechanisms against development-set exploitation, and failed recursive generations are often more informative than an uninterrupted sequence of wins.

The V20 result also gives a concrete example of why process utility deserves separate measurement. In many expensive discovery systems, knowing when \emph{not} to evaluate is itself valuable. G2's advantage is abstention: on A016 it stops because the observed state already exceeds its learned threshold. The control policy continues, requests two plausible mutations, and the hidden evaluator confirms that neither improves the champion. Under the frozen utility, the abstaining policy is therefore better.

Whether this kind of efficiency recursion can lead to recursive quality uplift is unknown. One possibility is that a better search policy first saves budget on saturated regions and later reallocates that budget to higher-value branches. Dream-RSI reports related adaptive changes in exploration breadth and cost over recursive rounds \cite{zheng2026dreamrsi}. Demonstrating that stronger effect in Genesis would require a new prospective protocol in which later generations actually discover higher-quality real-organism champions under an equal or otherwise controlled budget.

The next scientifically meaningful targets are:
\begin{enumerate}
  \item repeat the G1$\rightarrow$G2 comparison from multiple fresh real campaign states;
  \item test the frozen G2 policy under a new evaluator/task family;
  \item automate and attest proposer isolation;
  \item define an equal-budget recursive-quality gate separate from \lthreep;
  \item attempt G2$\rightarrow$G3 only under a protocol that cannot redefine the already established claim.
\end{enumerate}

\paragraph{Conclusion.}
Genesis IV establishes a prospectively validated, bounded two-transition recursive \emph{process}-improvement chain. G0$\rightarrow$G1 reduces hidden-evaluation cost from four to two observations at unchanged retained quality. A lineage-mediated G2 then survives a fresh 24-world replay holdout with 24 paired wins and, on the real A016 state, matches G1's retained outcome while requesting zero hidden evaluations rather than two. This satisfies the project's preregistered \lthreep definition. It does not establish open-ended RSI, general transfer, or recursive quality uplift.

\section*{AI-assistance disclosure}

Generative-AI systems were used extensively during software development, proposer execution, adversarial review, data inspection, and manuscript preparation. Anthony Mets remains responsible for research objectives, protocol decisions, acceptance or rejection of generated work, scientific claims, and publication. AI-assisted internal review is not represented as independent peer review or external replication.

\section*{Artifact availability}

The publication package is maintained at \url{https://github.com/mjodheim/mira-genesis}. The frozen experimental evidence for V14--V20 is preserved in the public \code{mjodheim/mira} repository on the \code{research/genesis-dream-replay} line, including preregistrations, source freezes, holdout records, immutable physical-evaluation branches, the evidence register, and post-result audit artifacts.
